\documentclass[UKenglish, thumbmarks]{uiophdthesis}
\usepackage[utf8]{inputenc}
\usepackage[T1]{fontenc, url}  \urlstyle{sf}
\usepackage{babel, csquotes, fancyvrb, graphicx, textcomp, uiomasterfp, varioref}
%% \usepackage[backend=biber,style=numeric-comp]{biblatex}  %% No bibliography
\usepackage[hidelinks, hypertexnames=false]{hyperref}

\newcommand{\bsl}{\textbackslash}
\newcommand{\dash}{\textthreequartersemdash}
\newcommand{\p}[1]{\textsf{#1}}
\newcommand{\pb}[1]{\textbf{\p{#1}}}
\newcommand{\pcmd}[1]{\p{\bsl #1}}
\newcommand{\popt}[1]{\p{[#1]}}
\newcommand{\ppar}[1]{\p{\{#1\}}}
\newcommand{\tighterlist}{\addtolength{\itemsep}{-1ex}}

\title{Writing your PhD thesis}
\subtitle{A guide to the \LaTeX{} document class \pb{uiophdthesis}}
\author{Dag Langmyhr (\url{dag@ifi.uio.no})}

\begin{document}
\uiomasterfp[kind=Documentation, compact, nosp, colour=grey,
  date=\today]

\frontmatter
\maketitle

\begin{abstract}
\LaTeX{} is an excellent tool for writing your PhD thesis, especially in
combination with the bibliography tool Bib\LaTeX.

The University of Oslo has defined guidelines\footnote{\label{guidelines}See
  \url{https://www.uio.no/english/services/print/doctoral-thesis/index.html}.}
for the PhD theses, and 
the University of Oslo Library and the Department of Informatics have
developed the document class \pb{uiophdthesis} in accordance with these
guidelines. This documentation has been typeset using that class.
\end{abstract}

\tableofcontents

\mainmatter
\begin{uioseparatorpage}{Introduction\label{intro-part}}
  What is the \pb{uiophdthesis} document class and how do I use it?
\end{uioseparatorpage}

\chapter{About the document class}
The \p{uiophdthesis} class has the following features:
\begin{itemize}
\item The paper size is A4, which is what the printers require (see
  the guidelines in footnote on page~\pageref{guidelines}).
  The thesis will be photographically reduced to
  81\,\% when printed.

\item The font size is 12\,pt which will look all right in the reduced
  print.

\item You should not attempt to create any kind of cover page. This
  will always be supplied by the printers.
\end{itemize}

\section{Document structure}
The standard \LaTeX{} document structuring commands
\begin{center}
  \pcmd{chapter}, \pcmd{section}, \pcmd{subsection} and \pcmd{subsubsection}
\end{center}
can be used in your thesis.

You should avoid using the \pcmd{part} command. If you want a
separation page, you should rather use the \p{uioseparatorpage}
environment:

\begin{quote}
  \pcmd{begin}\ppar{uioseparatorpage}\ppar{\emph{title}}\\
  ~~~~ some text (if you need it)\\
  \pcmd{end}\ppar{uioseparatorpage}
\end{quote}
For an example, see page~\pageref{intro-part}.

\section{Included research papers\label{uiopaper}}
Your thesis is likely to contain one or more research paper published
in various journals. When you want to include such a research
paper, you should use the command \textbf{\pcmd{uiopaper}}:
\begin{quote}
  \pcmd{uiopaper}\popt{\emph{options}}\ppar{\emph{title of paper}}
\end{quote}
The available options are:
\begin{description}
\item[\p{subtitle=\ppar{Any suitable subtitle}}] provides a subtitle
  (if required).
\item[\p{authors=\ppar{My Name\pcmd{and} Another Name\bsl\bsl\ affiliation}}] is
  used when you want to name the author or authors.
\item[\p{abstract=\ppar{This article describes \dots}}] provides
  an abstract for the article (if you want one on the paper title
  page).
\item[\p{compact}] gives a more compact title page~\dash{}
  this is particularly useful if you have a long
  abstract. Page~\pageref{info-part} has a default header page, while
  page~\pageref{part:example} uses the compact version.
\item[\p{published=\ppar{Published in \emph{Journal of ABC}\bsl\bsl\ 
      additional info}}] for any information on the publication of the
  article. 
\item[\p{info=\ppar{\dots}}] for any additional information.
\item[\p{nochapter}] see Section~\vref{article-style}.
\end{description}

\subsection{Typesetting a research paper}
If you have the \LaTeX{} source for the research paper, you can include it
here, after the \pcmd{uiopaper} command.

\subsubsection{Research papers in \emph{article} style\label{article-style}}
Many research papers were originally written in the \p{article} style,
and this style has no \pcmd{chapter} command. To make it easier to include
these in the thesis, the \pcmd{uiopaper} command may be given as
\begin{quote}
  \p{\pcmd{uiopaper}[nochapter, \dots]\ppar{\emph{paper title}}}
\end{quote}
Then, the paper will be properly typeset with no \pcmd{chapter}
commands but with \pcmd{section} as the highest structuring command.

\subsection{Research paper published elsewhere}
If you want to include the research paper exactly as it was typeset in
some publication, you can import it from a \textsc{pdf} file using the
command \textbf{\pcmd{uioincludepdf}} giving the \textsc{pdf} file as
parameter.

Page numbers can often be a problem when including full-page
\textsc{pdf} files: they might overlap those of your thesis. To remedy
this, you may supply \pcmd{uioincludepdf} with options to alter the
position of page number of the thesis. (Those of the included paper are, of
course, part of the \textsc{pdf} file and cannot be modified.)

\medskip\noindent
\p{\pcmd{uioincludepdf}\ppar{\emph{file name}}} --- move thesis
  page numbers into the margin\\
\p{\pcmd{uioincludepdf}[numbers=low]\ppar{\emph{file name}}} --- as
  above, but lower\\
\p{\pcmd{uioincludepdf}[numbers=none]\ppar{\emph{file name}}} ---
  remove thesis page numbers

\uiopaper[authors={My Name\and Your Name\\ The University of Oslo},
    abstract={This ``paper'' provides instructions on how to use this
      document class. It also describes the available document options.},
    info={(This page is an example of a default header page for
      published papers.)}]
  {Practical information\label{info-part}}
  
\chapter{Installation}
\section{On a Linux computer at Ifi}
If you are processing your \LaTeX{} document on a stationary Linux
computer at the University of Oslo, you need not worry about
installing the \p{uiophdthesis} document class; it is already there.

\section{On your personal computer\label{privat-pc}}
To use this document class on your own computer (which may run Linux, MacOS
or Windows) you must do the following:
\begin{enumerate}
\item Fetch
  \url{https://www.mn.uio.no/ifi/tjenester/it/hjelp/latex/uiotheses.zip}. (Click
  on this \textsc{url} to download the file.)
  
\item Unzip the files. You may place all the files in the same folder
  as your \LaTeX{} source files.\footnote{If you know where \LaTeX{}
    packages are kept on your computer, you 
    can save them there to make them generally available. Remember to
    refresh your file name database afterwards.}
\end{enumerate}
And that should be all.

\section{Using Overleaf}
If you are using Overleaf (see \url{https://www.overleaf.com}) to
write your thesis, you must do the following:
\begin{enumerate}
\item Select ``New Project''.

\item Under ``Institution Templates'', select ``University of Oslo''.

\item Select ``UiO PhD thesis''.

\item Select ``Open as Template''.
\end{enumerate}
Once this has been done, you may use the document class.

\chapter{Using the document class}
To use this document class, just start your \LaTeX{} file with
\begin{quote}
  \pcmd{documentclass[\emph{options}]}\ppar{uiophdthesis}
\end{quote}

\section{Document class options}

\subsection{Fonts\label{font}}
The default document font for \p{uiophdthesis} is the standard \LaTeX{} font
{\fontfamily{lmr}\fontshape{it}\selectfont Computer Modern} in
combination with \textsf{\textit{Helvetica}}.
You may select a different font by using one of these
options to \pcmd{documentclass}:
\begin{description}\tighterlist
\item[\p{font=garamond}] selects
  {\fontfamily{mdugm}\fontshape{it}\selectfont Garamond} with
  \textsf{\textit{Helvetica}}
  
\item[\p{font=noto}] selects
  {\fontfamily{NotoSerif-LF}\fontshape{it}\selectfont Noto serif} with
  {\fontfamily{NotoSans-TLF}\fontshape{sl}\selectfont Noto sans}

\item[\p{font=times}] selects
  {\fontfamily{ptm}\fontshape{it}\selectfont Times Roman} with
  \textsf{\textit{Helvetica}}
\end{description}

\subsection{Option \p{thumbmarks}\label{lab:thumbmarks}}
This option provides highly visible thumb marks on the front page for
all papers included using the command \pcmd{uiopaper}
(see Section~\vref{uiopaper}). Pages~\pageref{info-part} and
\pageref{part:example} show such thumb marks.

\subsection{Other document class options}
All other options to \pcmd{documentclass} are passed on to
the packages you use.

\uiopaper[compact,
    author={A Name\\ University of Somewhere},
    abstract={This ``paper'' shows a suitable base text for a
      \LaTeX{} file for a PhD thesis, and it describes the various
      specifications used, line for line.},
    info={(This page is an example of a compact header page
      for published papers; see Section~\vref{lab:thumbmarks}.)}]
  {An example\label{part:example}}
\noindent
The \p{uiophdthesis} package comes with a base file named
\pb{uiophdthesis-base.tex} containing the basic layout of your thesis;
see Figures~\vref{fig:base} and~\vref{fig:base2}.
  The idea is that you make a copy of that
  file, modify the specified texts, and then write your thesis.

\begin{description}
\item[Line 1:] The document class should be \pb{uiophdthesis}. You must
  also specify the language of your thesis. Include option \pb{thumbmarks} if
  you want thumbmarks; see Section~\vref{lab:thumbmarks}.  
\item[Line 2:] UTF-8 is the most common character encoding in use
  today, so, unless you specify otherwise in your text editor, you are
  likely to get this encoding.
\item[Line 3:] The \p{url} package provides the \pcmd{url}
  command which is very useful for typesetting long internet
  addresses. These should be set in a \textsf{sans serif} (``sf'') typeface
  (rather than \texttt{teletype}). For an example, see
  Section~\vref{privat-pc}.
\item[Lines 4--8:] These packages should always be included:
  \begin{description}
  \item[\p{babel}] handles language adaption.
  \item[\p{csquotes}] supports quote marks in various languages. This
    package is required by \p{biblatex}; see below.
  \item[\p{graphicx}] provides support for including illustrations.
  \item[\p{textcomp}] adds many useful symbols.
  \item[\p{varioref}] gives improved features for crossrefererencing.
  \item[\p{biblatex}] loads Bib\LaTeX{} which handles
    bibliographies.\footnote{\emph{Local guide to Bib\LaTeX} at
      \url{https://www.mn.uio.no/ifi/tjenester/it/hjelp/latex/biblatex-guide.pdf}
      is a simple introduction to creating your bibliography.}
    The package options given here are recommended; they use the
    numeric citation style favoured in natural science.
  \item[\p{hyperref}] provides hyperlinks both internally and externally.
  \item[\p{bookmark}] fixes some minor bugs in \p{hyperref}.
  \end{description}

\item[Lines 10--11:] You must always state a thesis title. You may also
  provide a subtitle, but that is not mandatory.
\item[Line 12:] Don't forget you own name! If there are several
  authors, separate the names with an \pcmd{and}.

\item [Line 14:] \pcmd{addbibresouce} specifies the name/s of your Bib\LaTeX{}
  bibliography file/s.

\item[Line 17:] \pcmd{frontmatter} specifies the start of the thesis front matter, i.e.,
  abstract, table of contents~etc.
\item[Line 18:] prints a title page. (This will appear after
  the cover page so it is recommended.) Normally, just the title, subtitle and
  author's name is printed, but you may use these options to
  provide additional information:
  \begin{description}
  \item[\p{dept=\ppar{\emph{department name}}}] gives the name of the
    department. 
  \item[\p{fac=\ppar{\emph{faculty name}}}] gives the faculty name. 
  \item[\p{info=\ppar{\emph{additional info}}}] provides additional information.
  \item[\p{supervisor=\ppar{\emph{supervisor's name}}}] names the
    supervisor. 
  \item[\p{supervisors=\ppar{\emph{supervisor's name}\pcmd{and}
        \emph{name}\pcmd{and} \emph{name}}}] gives the names of the
    supervisors, if there are more than one.
  \item[\p{year=\ppar{\emph{year}}}] indicates the year; it is only
    required when the year is other than the current one.
  \end{description}

\item[Lines 27--29:] contains your abstract.
\item[Lines 30--33:] contains your abstract in a different language,
  for example in Norwegian. \textbf{A Norwegian abstract is mandatory
    at the University of Oslo.} 
\item[Lines 35--37:] produces your tables of content, figures and
  tables, accordingly.
\item[Lines 39--41:] is your preface.

\item[Line 43:] shows the start of the main part of your thesis.
\item[Line 44--46:] gives you a separation page.
\item[Line 48--49:] are the standard structuring commands.
\item[Line 51--59:] contains a research paper typeset as part of the thesis.
\item[Line 61--62:] contains a different research paper avaible as a
  \textsc{pdf} file. 
\item[Line 64--67:] gives another separator page followed by some
  text.
\item[Line 71:] starts the back part containing appendices,
  bibliography and such.
\item[Line 72:] prints the bibliography created by Bib\LaTeX.
\end{description}

\begin{figure}[htp]
  \VerbatimInput[fontsize=\footnotesize,frame=single,obeytabs,
    label={\pb{uiophdthesis-base.tex}},numbers=left,lastline=41]{uiophdthesis-base.tex}
  \caption{The file \p{uiophdthesis-base.tex} (part I)\label{fig:base}}
\end{figure}

\begin{figure}[htp]
  \VerbatimInput[fontsize=\footnotesize,frame=single,obeytabs,
    numbers=left,firstline=43]{uiophdthesis-base.tex}
  \caption{The file \p{uiophdthesis-base.tex} (part II)\label{fig:base2}}
\end{figure}

\chapter{Two other examples}
\begin{itemize}
\item The file \url{uiophdthesis-guide.tex} shows the \LaTeX{} source
code for this documentation.
\item The file \url{uiophdthesis-demo.pdf} presents a demonstration
  thesis showing both the \LaTeX{} code and the result.
\end{itemize}
\end{document}
